\chapter{Introduction}

\section{Motivation: attention is the scarce resource}
In the Mah\=abh\=arata, Dro\d{n}a, the martial teacher of the P\=a\d{n}\d{d}ava and Kaurava princes, sets up a wooden 
bird in a tree and tells his students to aim for its eye. Before each releases the arrow, he asks what 
they see. One says the tree. Another the branches, the leaves, his brothers standing beside him. 
Dro\d{n}a asks Arjuna the same question, then asks again: does he see the tree? The branches? Arjuna answers
no to each, until only one thing remains: "I see only the eye of the bird." For that moment, the tree,
the branches, and everyone else standing there have not merely fallen out of Arjuna's interest. They have gone.
This dissertation pursues the same operation on a headset instead of a bow: not adding emphasis to what matters,
 but removing everything that does not, until it is no longer competing to be seen.

The modern visual environment is engineered against sustained attention. Open-plan offices place
moving colleagues at the edge of vision. Shared homes put televisions and phone screens beside
workspaces. Streetscapes compete for a driver's gaze with advertising designed by professionals to
win it. The cost is well documented: task-irrelevant peripheral events capture attention
involuntarily, and each capture carries a performance penalty on the task the viewer intended to do.
In immersive settings the effect is directly measurable, introducing peripheral visual distractors
into a virtual classroom more than doubled commission errors on a sustained-attention task in a
recent study of sixty-six participants \citep{ai2025}, replicating decades
of laboratory findings on involuntary attentional capture by peripheral motion and luminance change
\citep{itti1998}.

The conventional response of the extended-reality (XR) community to attention problems has been
\textit{additive}: arrows, halos, attention funnels, and highlights layered onto the world to point at what
matters \citep{biocca2006}. This dissertation pursues the complementary and far less explored
operation: \textit{subtraction}. If the periphery pulls attention away, dim it, mute it, or re-grade it so
that it pulls less. In the literature this family of operations is called diminished reality (DR), 
concealing or attenuating real-world content rather than adding to it \citep{mori2017,cheng2022},  and its application to distraction has recently been named \textit{visual noise
cancellation}, by direct analogy with the acoustic kind \citep{hong2024}.

A video-see-through (VST) headset is, in principle, the ideal instrument for visual noise
cancellation. Unlike optical see-through (OST) glasses, whose additive light engines physically
cannot darken the world, a VST headset re-displays the world on an opaque screen: every pixel the
user sees is, in principle, available for modification. The Meta Quest 3, an affordable, widely
owned consumer device with colour passthrough, would seem to make DR-based attention guidance
deployable at scale for the first time. In principle.

\section{The compositing constraint}

In practice, the premise "every pixel is available for modification" is false on consumer hardware,
and that falsity is the motivating technical problem of this dissertation.

On the Quest 3, the passthrough view of the world is composited by the operating system, not by the
application. An application cannot read the passthrough layer, cannot shade it, and cannot spatially
mask it. It may only draw \textit{on top} of it, or apply a global colour transform through a narrow styling
interface. Since early 2025 the platform has additionally exposed the raw forward camera feed to
applications (Passthrough Camera Access), but at materially lower quality than the system passthrough
the user normally sees: a monocular 1280$\times$960 stream covering roughly 85–90$^{\circ}$ of field of view, versus
the system layer's full coverage of the headset's approximately 110$^{\circ}$ display, with its correct
per-eye reprojection. The one place where full pixel control exists is therefore also the place where
image quality is worst. A recent psychophysical study shows the stakes plainly: no current VST
headset reaches normal human visual acuity through its passthrough, with the Quest 3 degrading
further in low light \citep{wang2026}.

Subtractive attention guidance on this platform is consequently a \textit{compositing-rights problem}. The
effect one wants: "dim everything except the region that matters, without degrading the region that
matters", cannot be obtained by editing the image the user sees, because no application is permitted
to edit it. Chapter 3 shows that what remains is a narrow but genuinely usable design space, organised
into three tiers of access, and Chapter 4 presents a working system that exploits an architectural
inversion at its centre: the region that matters is rendered as an \textit{absence},  a transparent window
in an overlay, so that the highest-quality image on the device, the native system passthrough,
serves the focus region for free, while the manipulation is confined to the periphery where acuity is
lower and precision matters less. On the same principle, the system's most capable mode composites
native passthrough inside the window with a shader-processed camera feed outside it, a hybrid for
which a structured prior-art search (Chapter 2) found no published or shipped precedent.

\section{One question, two halves}

The dissertation asks a single question, with two distinct halves:

\begin{quote}
\textbf{Can subtractive attention guidance be delivered on consumer XR hardware? And when it is
delivered, does it actually help human attention?}
\end{quote}

The two halves of the question structure the entire work.

\textbf{Part T (technical): can it be built?} What degree of diminishment control is achievable on a
consumer VST headset without OS-level access to the passthrough layer, and at what cost to image
quality, field of view, latency, and comfort? Part T is answered by the design space (Chapter 3), the
reference implementation (Chapter 4), and the technical evaluation plan with preliminary results
(Chapter 5).

\textbf{Part H (human): does it work?} Under controlled distraction, does peripheral diminishment reduce
the cost distractors impose on a focal task, and does salience re-grading enhance focal detection
without unacceptably degrading peripheral awareness? Part H is answered by a pre-registered user
study design (Chapter 6).

Two methodological commitments connect the halves and should be named at the outset, because they are
choices rather than accidents.

First, the study's dynamic scenario runs as a \textit{simulation}, pre-recorded driving footage
presented on the system's rendering sphere rather than live passthrough of a street, following the
closest published DR evaluations, which simulated diminishment inside virtual environments because
controlled, repeatable, event-dense stimuli cannot be obtained from the uninstrumented real world
\citep{cheng2022,mclaughlin2025}. Second, detection is supplied by an \textit{oracle}
computed offline over the study footage rather than by a model running on the headset, so the study
can ask whether an effect built on excellent detection helps attention without first engineering
excellent detection on-device. Chapter 5 measures how far a live pipeline currently falls short of
that oracle, and Chapter 6 carries both commitments forward as stated limitations.

The two study blocks bracket the feasibility spectrum. The primary, workstation, block runs on
\textit{real passthrough with real physical distractors}, the actual artefact with no simulation. The
secondary, driving, block evaluates the concept in the simulated deployment context.

\section{Research questions and hypotheses}

Three research questions operationalise the spine question. They are stated here in summary. Their
full pre-registered form, including smallest effect sizes of interest, equivalence margins,
statistical tests, and a decision table committing the dissertation to a conclusion under every
outcome, is given in Chapter 6.

\begin{itemize}
\item \textbf{RQ1 (primary).} Does peripheral DR dimming (Hard Dark, world-locked window) reduce the
  performance cost that peripheral visual distractors impose on a focal workstation task?
  \textit{Hypothesis H1:} the error rate on the workstation task, run with distractors present
  throughout, is lower with the vignette on than off, by at least a pre-registered fraction of the
  vignette-off error rate.
\item \textbf{RQ2 (secondary).} Does SignPop's detection-gated salience re-grading improve detection
  specifically where the gate has room to act, without degrading detection pooled across the whole
  scene beyond an acceptable margin?
  \textit{Hypotheses H2a/H2b:} detection improves in the eccentricity band where the colour-pop gate is
  strong but the target is still within useful glance range (H2a), and hit rate pooled across the
  whole field of view is equivalent to the vignette-off condition within a 10-percentage-point margin
  (H2b). H2b is a safety hypothesis tested by equivalence, so failing it would be a conclusive
  negative finding rather than an ambiguous null.
\item \textbf{RQ3 (exploratory).} After experiencing the system, what do participants report about
  perceived focus benefit, perceived awareness cost, and visual comfort?
\end{itemize}

The design philosophy behind these hypotheses deserves one sentence here: the study is built so that
it cannot return "inconclusive". Distractors are already known to impose a measurable cost, on the
evidence of the sixty-six-participant result cited above. The study takes that cost as given, asks
whether the system reduces it, and pre-commits to the interpretation of every outcome, including
failure.

\section{Contributions}

\begin{itemize}
\item \textbf{C1: Design space.} A taxonomy of DR capability under compositing constraints on consumer VST
  hardware: three access tiers, with the achievable effects, ceilings, and costs of each,
  characterised empirically rather than speculatively (Chapter 3).
\item \textbf{C2: Reference implementation.} An open, working multi-mode DR attention-guidance system on
  Quest 3 passthrough: the world-locked focus-window-as-absence architecture, world-direction fusion
  sampling that lets a monocular camera feed fuse binocularly, motion-coupled comfort suppression,
  and ten peripheral treatments spanning the three-tier space. Two of them, SignPop (detection-gated
  salience re-grading) and Hard Dark (full peripheral occlusion), are carried to confirmatory
  evaluation, with documented failure modes (Chapter 4).
\item \textbf{C3: Technical evaluation.} A benchmark plan (frame cost, latency, legibility across rendering
  paths, thermal endurance, and the oracle-versus-live detection gap) with preliminary results from
  the offline detection pipeline (Chapter 5).
\item \textbf{C4: Pre-registered user study, with interim results.} A falsifiable within-subjects design
  (N = 20) with a single properly powered primary hypothesis, equivalence-bounded safety testing,
  pilot gates, and a pre-committed decision table, together with an interim analysis of the fourteen
  participant pairs collected so far: a workstation-focus benefit (+3.34 points accuracy), no loss of
  peripheral awareness, and a detection gain specifically in the eccentricity band the system's own
  design predicts (Chapter 6).
\end{itemize}

\section{Scope: what this dissertation does not claim}

Because the credibility of a small, carefully scoped study depends on the discipline of its claims,
the non-claims are stated as prominently as the claims.

\textbf{No driving or road-safety claims.} The driving footage used in the secondary study block is a
dynamic, high-event-rate monitoring \textit{stimulus}. A seated participant passively watching
passenger-perspective video licenses no inference about driving, drivers, or road safety. Such
inference would require a driving simulator with vehicle control, which is out of scope.

\textbf{No product claims.} The study evaluates attention mechanisms under controlled distraction, not the
desirability or readiness of a wearable product. Comfort findings from a 2026 headset transfer to
future lighter hardware only as bounds, and are reported as such.

\textbf{Oracle conditioning stays visible.} SignPop's finding depends on detection: on offline,
oracle-quality detection, not on live, on-device detection. Every place this dissertation reasons
about a live or on-device detector, that reasoning is conditioned on the oracle bound that Chapter 5
measures, and marked accordingly. Oracle performance is never substituted for live performance
without saying so.

\textbf{Interim, not final, results.} At the time of writing the technical benchmark campaign of
Chapter 5 has not been run, beyond the offline detection bake reported there as a preliminary result.
The user study of Chapter 6 is a different matter. It is actively collecting, against the
pre-registered design that chapter specifies, and eleven usable pairs exist for each of its two
confirmatory blocks against a target of twenty. Chapter 6, Section 6.11 reports what those eleven
pairs show, explicitly as an interim analysis. Recruitment continues, achieved power at this $N$ is
below the pre-registered target, and every claim built on it in this dissertation is qualified as
interim accordingly. No result in this document is asserted as the pre-registered confirmatory
finding. Where an interim result is discussed, it is named as interim every time.

\section{Roadmap}

Chapter 2 surveys the literature across diminished reality, saliency modulation and gaze guidance,
foveated rendering, attention guidance in driving and workstation contexts, and human-factors
foundations, closing with the research gaps this work addresses and a documented prior-art search for
the system's core architecture. Chapter 3 develops the three-tier design space of DR under
compositing constraints. Chapter 4 presents the system: architecture, the full mode family and the
three modes the study uses, interaction design, and the engineering record including dead ends. Chapter 5 specifies the technical evaluation
and reports preliminary detection-pipeline results. Chapter 6 presents the pre-registered user study
design in full and reports interim results from the fourteen Block A and thirteen Block B pairs
collected so far.
Chapter 7 discusses these interim results, threats to validity, and transfer, including what does
and does not carry to optical see-through hardware. Chapter 8 concludes.

\figplaceholder{60mm}{The three peripheral treatments seen from the user's viewpoint: SignPop, Soft Dark, Hard Dark, each with the world-locked focus window visible. Composite of headset screenshots.}{The three peripheral treatments seen from the user's viewpoint.}

\figplaceholder{60mm}{The two-part structure of the dissertation: Part T (design space $\rightarrow$ system $\rightarrow$ benchmarks) and Part H (pre-registered study), with the workstation block bridging them on real hardware.}{The two-part structure of the dissertation.}
